
\[
\begin{tikzcd}[column sep=large,row sep=large]
D_+(st) \arrow[rr,hook,"i"] \arrow[dr,"f_{U,st}"'] & &
D_+(s) \arrow[dl,"f_{U,s}"]\\
& U &
\end{tikzcd}
\]
il s'%
agit de montrer que les Modules \(i^*\circ\bar f_{U,s}^{\,*}\bigl(\underline{M}_{(s),U}\bigr)\) et \(\bar f_{U,st}^{\,*}\bigl(\underline{M}_{(st),U}\bigr)\)
sont canoniquement isomorphes. Or on a d'une part les isomorphismes canoniques :
\[
\begin{aligned}
i^*\circ\bar f_{U,s}^{\,*}\bigl(\underline{M}_{(s),U}\bigr)
&\simeq
i^*\Bigl(f_{U,s}^{-1}\bigl(\underline{M}_{(s),U}\bigr)
 \otimes_{f_{U,s}^{-1}((B|U)_{(s)})}\underline{O}_{D_+(s)}\Bigr)\\
&\simeq
i^{-1}\Bigl(i_{U,s}^{-1}\bigl(\underline{M}_{(s),U}\bigr)
 \otimes_{f_{U,s}^{-1}((B|U)_{(s)})}\underline{O}_{D_+(s)}\Bigr)
 \otimes_{i^{-1}(\underline{O}_{D_+(s)})}\underline{O}_{D_+(st)}\\
&\simeq
f_{U,st}^{-1}\bigl(\underline{M}_{(s),U}\bigr)
 \otimes_{f_{U,st}^{-1}((\underline{B}|U)_{(s)})}\underline{O}_{D_+(st)} .
\end{aligned}
\]
D'autre part :
\[
  \bar f_{U,st}^{\,*}\bigl(\underline{M}_{(st),U}\bigr)
  \simeq
  f_{U,st}^{-1}\bigl(\underline{M}_{(st),U}\bigr)
  \otimes_{f_{U,st}^{-1}((\underline{B}|U)_{(st)})}\underline{O}_{D_+(st)} .
\]
Enfin, il résulte de EGA II 2.2.2, traduit en termes de faisceaux, qu'on a les isomorphismes canoniques :
\[
  \underline{M}_{(st),U}\simeq
  \bigl(\underline{M}_{(s),U}\bigr)_{t^d/s^e}
  \simeq
  \underline{M}_{(s),U}\otimes_{(\underline{B}|U)_{(s)}}(\underline{B}|U)_{(st)},
\]
d'où finalement l'isomorphisme cherché.

Il est clair que si \(\underline{M}\) est quasi-cohérent le Module \(\Projz(\underline{M})\) ainsi construit est bien celui qui est défini dans EGA II.

2.3.3. Il résulte de 2.3.1 et 2.3.2 que le foncteur \(\Projz\) est exact. Revenant aux notations de 2.2, on en déduit donc par passage aux catégories dérivées un foncteur encore noté \(\Projz\) :
\[
  \Projz:D^+(\mathcal C)\longrightarrow D^+(\mathcal A).
\]

On voit également sur les constructions précédentes que le foncteur
